% problem_id: S001-artin-lemma-limit-preserving
% source_id: S001 (Stacks Project)
% source_file: artin.tex
% statement_locator: artin.tex lines 5418--5423
% proof_locator: artin.tex lines 5425--5955
% env: lemma
% id_note: label 跨文件重名 ⇒ id 用文件限定形式
% pairing: tight (verified)
% status: complete (statement + author proof verbatim + reason + locators)
%
\begin{lemma}
\label{lemma-limit-preserving}
In Situation \ref{situation-contractions} the functor $F$ is limit preserving:
for any directed limit $V = \lim V_\lambda$ of Noetherian affine schemes
over $S$ we have $F(V) = \colim F(V_\lambda)$.
\end{lemma}

\begin{proof}
This is an absurdly long proof. Much of it consists of standard
arguments on limits and \'etale localization. We urge the reader to skip ahead
to the last part of the proof where something interesting happens.

\medskip\noindent
Let $V = \lim_{\lambda \in \Lambda} V_i$ be a directed limit of schemes
over $S$ with $V$ and $V_\lambda$ Noetherian and with affine transition
morphisms. See Limits, Section \ref{limits-section-limits} for material on
limits of schemes. We will prove that $\colim F(V_\lambda) \to F(V)$
is bijective.

\medskip\noindent
Proof of injectivity: notation.
Let $\lambda \in \Lambda$ and
$\xi_{\lambda, 1}, \xi_{\lambda, 2} \in F(V_\lambda)$ 
be elements which restrict to the same element of $F(V)$.
Write $\xi_{\lambda, 1} = (Z_{\lambda, 1}, u'_{\lambda, 1},
\hat x_{\lambda, 1})$ and
$\xi_{\lambda, 2} = (Z_{\lambda, 2}, u'_{\lambda, 2}, \hat x_{\lambda, 2})$.

\medskip\noindent
Proof of injectivity: agreement of $Z_{\lambda, i}$.
Since $Z_{\lambda, 1}$ and $Z_{\lambda, 2}$ restrict to the same closed
subset of $V$, we may after increasing $i$ assume
$Z_{\lambda, 1} = Z_{\lambda, 2}$, see
Limits, Lemma \ref{limits-lemma-inverse-limit-top} and
Topology, Lemma \ref{topology-lemma-describe-limits}.
Let us denote the common value $Z_\lambda \subset V_\lambda$, for
$\mu \geq \lambda$
denote $Z_\mu \subset V_\mu$ the inverse image in $V_\mu$ and
and denote $Z$ the inverse image in $V$. We will use below that
$Z = \lim_{\mu \geq \lambda} Z_\mu$ as schemes if we view $Z$
and $Z_\mu$ as reduced closed subschemes.

\medskip\noindent
Proof of injectivity: agreement of $u'_{\lambda, i}$.
Since $U'$ is locally of finite type over $S$ and since
the restrictions of $u'_{\lambda, 1}$ and $u'_{\lambda, 2}$
to $V \setminus Z$ are the same, we may after increasing $\lambda$ assume
$u'_{\lambda, 1} = u'_{\lambda, 2}$, see Limits, Proposition
\ref{limits-proposition-characterize-locally-finite-presentation}.
Let us denote the common value $u'_\lambda$ and denote
$u'$ the restriction to $V \setminus Z$.

\medskip\noindent
Proof of injectivity: restatement.
At this point we have
$\xi_{\lambda, 1} = (Z_\lambda, u'_\lambda, \hat x_{\lambda, 1})$ and
$\xi_{\lambda, 2} = (Z_\lambda, u'_\lambda, \hat x_{\lambda, 2})$.
The main problem we face in this part of the proof
is to show that the morphisms
$\hat x_{\lambda, 1}$ and $\hat x_{\lambda, 2}$ become the same after
increasing $\lambda$.

\medskip\noindent
Proof of injectivity: agreement of $\hat x_{\lambda, i}|_{Z_\lambda}$.
Consider the morpisms
$\hat x_{\lambda, 1}|_{Z_\lambda}, \hat x_{\lambda, 2}|_{Z_\lambda} :
Z_\lambda \to W_{red}$.
These morphisms restrict to the same morphism $Z \to W_{red}$.
Since $W_{red}$ is a scheme locally of finite type over $S$
we see using Limits, Proposition
\ref{limits-proposition-characterize-locally-finite-presentation}
that after replacing $\lambda$ by a bigger index
we may assume
$\hat x_{\lambda, 1}|_{Z_\lambda} = \hat x_{\lambda, 2}|_{Z_\lambda} :
Z_\lambda \to W_{red}$.

\medskip\noindent
Proof of injectivity: end.
Next, we are going to apply the discussion in
Remark \ref{remark-diagonal} to $V_\lambda$ and the two elements
$\xi_{\lambda, 1}, \xi_{\lambda, 2} \in F(V_\lambda)$.
This gives us
\begin{enumerate}
\item $e_\lambda : E_\lambda' \to V_\lambda$
separated and locally of finite type,
\item $e_\lambda^{-1}(V_\lambda \setminus Z_\lambda) \to
V_\lambda \setminus Z_\lambda$ is an isomorphism,
\item a monomorphism $E_{W, \lambda} \to V_{\lambda, /Z_\lambda}$
which is the equalizer of $\hat x_{\lambda, 1}$ and $\hat x_{\lambda, 2}$,
\item a formal modification $E'_{\lambda, /Z_\lambda} \to E_{W, \lambda}$
\end{enumerate}
Assertion (2) holds by assertion (2) in Remark \ref{remark-diagonal}
and the preparatory work we did above getting
$u'_{\lambda, 1} = u'_{\lambda, 2} = u'_\lambda$.
Since $Z_\lambda = (V_{\lambda, /Z_\lambda})_{red}$ factors through
$E_{W, \lambda}$ because
$\hat x_{\lambda, 1}|_{Z_\lambda} = \hat x_{\lambda, 2}|_{Z_\lambda}$
we see from
Formal Spaces, Lemma \ref{formal-spaces-lemma-monomorphism-iso-over-red}
that $E_{W, \lambda} \to V_{\lambda, /Z_\lambda}$ is a closed immersion.
Then we see from assertion (4) in Remark \ref{remark-diagonal}
and Lemma \ref{lemma-closed-immersion} applied to the triple
$E_\lambda'$, $e_\lambda^{-1}(Z_\lambda)$,
$E'_{\lambda, /Z_\lambda} \to E_{W, \lambda}$ over $V_\lambda$ that
there exists a closed immersion $E_\lambda \to V_\lambda$
which is a solution for this triple.
Next we use assertion (5) in Remark \ref{remark-diagonal}
which combined with Lemma \ref{lemma-solution}
says that $E_\lambda$ is the ``equalizer'' of $\xi_{\lambda, 1}$
and $\xi_{\lambda, 2}$. In particular, we see that $V \to V_\lambda$
factors through $E_\lambda$. Then using Limits, Proposition
\ref{limits-proposition-characterize-locally-finite-presentation}
once more we find $\mu \geq \lambda$ such that $V_\mu \to V_\lambda$
factors through $E_\lambda$ and hence the pullbacks of
$\xi_{\lambda, 1}$ and $\xi_{\lambda, 2}$ to $V_\mu$ are the same
as desired.

\medskip\noindent
Proof of surjectivity: statement.
Let $\xi = (Z, u', \hat x)$ be an element of $F(V)$.
We have to find a $\lambda \in \Lambda$ and an element
$\xi_\lambda \in F(V_\lambda)$ restricting to $\xi$.

\medskip\noindent
Proof of surjectivity: the question is \'etale local.
By the unicity proved in the previous part of the proof and by the
sheaf property of $F$ in Lemma \ref{lemma-sheaf}, the problem
is local on $V$ in the \'etale topology. More precisely, let $v \in V$.
We claim it suffices to find an \'etale morphism
$(\tilde V, \tilde v) \to (V, v)$ and some
$\lambda$, some an \'etale morphism $\tilde V_\lambda \to V_\lambda$,
and some element $\tilde \xi_\lambda \in F(\tilde V_\lambda)$ such that
$\tilde V = \tilde V_\lambda \times_{V_\lambda} V$
and $\xi|_U = \tilde \xi_\lambda|_U$. We omit a detailed proof of
this claim\footnote{To prove this
one assembles a collection of the morphisms $\tilde V \to V$
into a finite \'etale covering and shows that the corresponding morphisms
$\tilde V_\lambda \to V_\lambda$ form an \'etale covering as well (after
increasing $\lambda$). Next one uses the injectivity to see that
the elements $\tilde \xi_\lambda$ glue (after increasing $\lambda$)
and one uses the sheaf property for $F$ to descend these
elements to an element of $F(V_\lambda)$.}.

\medskip\noindent
Proof of surjectivity: rephrasing the problem.
Recall that any \'etale morphism $(\tilde V, \tilde v) \to (V, v)$
with $\tilde V$ affine is the base change of an \'etale morphism
$\tilde V_\lambda \to V_\lambda$ with $\tilde V_\lambda$ affine
for some $\lambda$, see for example
Topologies, Lemma \ref{topologies-lemma-limit-fppf-topology}.
Given $\tilde V_\lambda$ we have
$\tilde V = \lim_{\mu \geq \lambda} \tilde V_\lambda \times_{V_\lambda} V_\mu$.
Hence given $(\tilde V, \tilde v) \to (V, v)$ \'etale with $\tilde V$ affine,
we may replace $(V, v)$ by $(\tilde V, \tilde v)$ and $\xi$ by
the restriction of $\xi$ to $\tilde V$.

\medskip\noindent
Proof of surjectivity: reduce to base being affine. In particular,
suppose $\tilde S \subset S$ is an affine open subscheme such
that $v \in V$ maps to a point of $\tilde S$. Then we may according
to the previous paragraph, replace $V$ by $\tilde V = \tilde S \times_S V$.
Of course, if we do this, it suffices to solve the problem
for the functor $F$ restricted to the category of locally Noetherian
schemes over $\tilde S$. This functor is of course the functor
associated to the whole situation base changed to $\tilde S$.
Thus we may and do assume $S = \Spec(R)$ is a Noetherian affine scheme
for the rest of the proof.

\medskip\noindent
Proof of surjectivity: easy case.
If $v \in V \setminus Z$, then we can take $\tilde V = V \setminus Z$.
This descends to an open subscheme $\tilde V_\lambda \subset V_\lambda$
for some $\lambda$ by Limits, Lemma
\ref{limits-lemma-descend-opens}.
Next, after increasing $\lambda$ we may assume
there is a morphism $u'_\lambda : \tilde V_\lambda \to U'$
restricting to $u'$. Taking
$\tilde \xi_\lambda = (\emptyset, u'_\lambda, \emptyset)$
gives the desired element of $F(\tilde V_\lambda)$.

\medskip\noindent
Proof of surjectivity: hard case and reduction to affine $W$.
The most difficult case comes from considering $v \in Z \subset V$.
We claim that we can reduce this to the case where $W$ is an
affine formal scheme; we urge the reader to skip this argument\footnote{Artin's
approach to the proof of this lemma is to work around this and
consequently he can avoid proving the injectivity first. Namely, Artin
consistently works with a finite affine \'etale coverings of all spaces
in sight keeping track of the maps between them during the proof.
In hindsight that might be preferable to what we do here.}.
Namely, we can choose an \'etale morphism
$\tilde W \to W$ where $\tilde W$ is an affine formal algebraic space
such that the image of $v$ by $\hat x : V_{/Z} \to W$ is
in the image of $\tilde W \to W$ (on reductions).
Then the morphisms
$$
p : \tilde W \times_{W, g} X'_{/T'} \longrightarrow X'_{/T'}
$$
and
$$
q : \tilde W \times_{W, \hat x} V_{/Z} \to V_{/Z}
$$
are \'etale morphisms of locally Noetherian formal algebraic spaces.
By (an easy case of) Algebraization of Formal Spaces, Theorem
\ref{restricted-theorem-dilatations-general}
there exists a morphism $\tilde X' \to X'$ of algebraic spaces
which is locally of finite type, is an isomorphism over $U'$, and
such that $\tilde X'_{/T'} \to X'_{/T'}$ is isomorphic to $p$.
By Algebraization of Formal Spaces, Lemma \ref{restricted-lemma-output-etale}
the morphism $\tilde X' \to X'$ is \'etale. Denote
$\tilde T' \subset |\tilde X'|$ the inverse image of $T'$.
Denote $\tilde U' \subset \tilde X'$ the complementary open subspace.
Denote $\tilde g' : \tilde X'_{/\tilde T'} \to \tilde W$
the formal modification which is the base change of $g$ by
$\tilde W \to W$. Then we see that
$$
\tilde X',\ \tilde T',\ \tilde U',\ \tilde W,
\ \tilde g : \tilde X'_{/\tilde T'} \to \tilde W
$$
is another example of Situation \ref{situation-contractions}.
Denote $\tilde F$ the functor constructed from this triple.
There is a transformation of functors
$$
\tilde F \longrightarrow F
$$
constructed using the morphisms $\tilde X' \to X'$ and
$\tilde W \to W$ in the obvious manner; details omitted.

\medskip\noindent
Proof of surjectivity: hard case and reduction to affine $W$, part 2.
By the same theorem as used above, there exists a morphism
$\tilde V \to V$ of algebraic spaces which is locally of finite type,
is an isomorphism over $V \setminus Z$ and such that
$\tilde V_{/Z} \to V_{/Z}$ is isomorphic to $q$.
Denote $\tilde Z \subset \tilde V$ the inverse image of $Z$.
By Algebraization of Formal Spaces, Lemmas
\ref{restricted-lemma-output-etale} and \ref{restricted-lemma-output-separated}
the morphism $\tilde V \to V$ is \'etale and separated.
In particular $\tilde V$ is a (locally Noetherian) scheme, see for example
Morphisms of Spaces, Proposition
\ref{spaces-morphisms-proposition-locally-quasi-finite-separated-over-scheme}.
We have the morphism $u'$ which we may view as a morphism
$$
\tilde u' : \tilde V \setminus \tilde Z \longrightarrow \tilde U'
$$
where $\tilde U' \subset \tilde X'$ is the open mapping isomorphically
to $U'$. We have a morphism
$$
\tilde {\hat x} :
\tilde V_{/\tilde Z} = \tilde W \times_{W, \hat x} V_{/Z}
\longrightarrow
\tilde W
$$
Namely, here we just use the projection. Thus we have the triple
$$
\tilde \xi = (\tilde Z, \tilde u', \tilde {\hat x}) \in \tilde F(\tilde V)
$$
We omit proving the compatibility condition; hints: if $V' \to V$, $\hat x'$,
and $x'$ witness the compatibility between $u'$ and $\hat x$, then one
sets $\tilde V' = V' \times_V \tilde V$ which comes with morphsms
$\tilde{\hat x}'$ and $\tilde x'$ and show this works.
The image of $\tilde \xi$ under the transformation $\tilde F \to F$
is the restriction of $\xi$ to $\tilde V$.

\medskip\noindent
Proof of surjectivity: hard case and reduction to affine $W$, part 3.
By our choice of $\tilde W \to W$, there is an affine open
$\tilde V_{open} \subset \tilde V$ (we're running out of notation)
whose image in $V$ contains our chosen point $v \in V$.
Now by the case studied in the next paragraph and the remarks made
earlier, we can descend $\tilde \xi|_{\tilde V_{open}}$
to some element $\tilde \xi_\lambda$ of $\tilde F$ over
$\tilde V_{\lambda, open}$
for some \'etale morphism $\tilde V_{\lambda, open} \to V_\lambda$
whose base change to $V$ is $\tilde V_{open}$.
Applying the transformation of functors $\tilde F \to F$
we obtain the element of $F(\tilde V_{\lambda, open})$
we were looking for. This reduces us to the case discussed
in the next paragraph.

\medskip\noindent
Proof of surjectivity: the case of an affine $W$. We have $v \in Z \subset V$
and $W$ is an affine formal algebraic space. Recall that
$$
\xi = (Z, u', \hat x) \in F(V)
$$
We may still replace $V$ by an \'etale neighbourhood of $v$.
In particular we may and do assume $V$ and $V_\lambda$ are affine.

\medskip\noindent
Proof of surjectivity: descending $Z$. We can find a $\lambda$
and a closed subscheme $Z_\lambda \subset V_\lambda$ such that
$Z$ is the base change of $Z_\lambda$ to $V$. See
Limits, Lemma \ref{limits-lemma-descend-finite-presentation}.
Warning: we don't know (and in general it won't be true)
that $Z_\lambda$ is a reduced closed subscheme of $V_\lambda$.
For $\mu \geq \lambda$ denote $Z_\mu \subset V_\mu$ the scheme theoretic
inverse image in $V_\mu$. We will use below that
$Z = \lim_{\mu \geq \lambda} Z_\mu$ as schemes.

\medskip\noindent
Proof of surjectivity: descending $u'$.
Since $U'$ is locally of finite type over $S$
we may assume after increasing $\lambda$
that there exists a morphism
$u'_\lambda : V_\lambda \setminus Z_\lambda \to U'$
whose restriction to $V \setminus Z$ is $u'$.
See Limits, Proposition
\ref{limits-proposition-characterize-locally-finite-presentation}.
For $\mu \geq \lambda$ we will denote $u'_\mu$ the restriction
of $u'_\lambda$ to $V_\mu \setminus Z_\mu$.

\medskip\noindent
Proof of surjectivity: descending a witness.
Let $V' \to V$, $\hat x'$, and $x'$ witness the compatibility between
$u'$ and $\hat x$. Using the same references as above we may assume
(after increasing $\lambda$) that there exists a morphism
$V'_\lambda \to V_\lambda$ of finite type whose base change to $V$
is $V' \to V$. After increasing $\lambda$ we may assume
$V'_\lambda \to V_\lambda$ is proper
(Limits, Lemma \ref{limits-lemma-eventually-proper}).
Next, we may assume $V'_\lambda \to V_\lambda$ is an isomorphism
over $V_\lambda \setminus Z_\lambda$
(Limits, Lemma \ref{limits-lemma-descend-isomorphism}).
Next, we may assume there is a morphism $x'_\lambda : V'_\lambda \to X'$
whose restriction to $V'$ is $x'$.
Increasing $\lambda$ again we may assume $x'_\lambda$
agrees with $u'_\lambda$ over $V_\lambda \setminus Z_\lambda$.
For $\mu \geq \lambda$ we denote
$V'_\mu$ and $x'_\mu$ the base change of $V'_\lambda$ and the
restriction of $x'_\lambda$.

\medskip\noindent
Proof of surjectivity: algebra.
Write $W = \text{Spf}(B)$, $V = \Spec(A)$, and
for $\mu \geq \lambda$ write $V_\mu = \Spec(A_\mu)$.
Denote $I_\mu \subset A_\mu$ and $I \subset A$
the ideals cutting out $Z_\mu$ and $Z$.
Then $I_\lambda A_\mu = I_\mu$ and $I_\lambda A = I$.
The morphism $\hat x$ determines and is determined by a
continuous ring map
$$
(\hat x)^\sharp : B \longrightarrow A^\wedge
$$
where $A^\wedge$ is the $I$-adic completion of $A$.
To finish the proof we have to show that this map descends to a map into
$A_\mu^\wedge$ for some sufficiently large $\mu$ where $A_\mu^\wedge$
is the $I_\mu$-adic completion of $A_\mu$.
This is a nontrivial fact; Artin writes in his paper
\cite{ArtinII}: ``Since the data (3.5) involve $I$-adic completions,
which do not commute with direct limits, the verification is somewhat
delicate. It is an algebraic analogue of a convergence proof in analysis.''

\medskip\noindent
Proof of surjectivity: algebra, more rings.
Let us denote
$$
C_\mu = \Gamma(V'_\mu, \mathcal{O})
\quad\text{and}\quad
C = \Gamma(V', \mathcal{O})
$$
Observe that $A \to C$ and $A_\mu \to C_\mu$ are finite ring maps as
$V' \to V$ and $V'_\mu \to V_\mu$ are proper morphisms, see
Cohomology of Spaces, Lemma
\ref{spaces-cohomology-lemma-proper-over-affine-cohomology-finite}.
Since $V = \lim V_\mu$ and $V' = \lim V'_\mu$
we have
$$
A = \colim A_\mu
\quad\text{and}\quad
C = \colim C_\mu
$$
by Limits, Lemma \ref{limits-lemma-descend-section}\footnote{We don't
know that $C_\mu = C_\lambda \otimes_{A_\lambda} A_\mu$ as the
various morphisms aren't flat.}. For an element
$a \in I$, resp.\ $a \in I_\mu$ the maps $A_a \to C_a$,
resp.\ $(A_\mu)_a \to (C_\mu)_a$ are isomorphisms by flat base change
(Cohomology of Spaces, Lemma
\ref{spaces-cohomology-lemma-flat-base-change-cohomology}).
Hence the kernel and cokernel of $A \to C$ is supported
on $V(I)$ and similarly for $A_\mu \to C_\mu$.
We conclude the kernel and cokernel of $A \to C$
are annihilated by a power of $I$ and the kernel and cokernel of
$A_\mu \to C_\mu$ are annihilated by a power of $I_\mu$, see
Algebra, Lemma \ref{algebra-lemma-Noetherian-power-ideal-kills-module}.

\medskip\noindent
Proof of surjectivity: algebra, more ring maps.
Denote $Z_n \subset V$ the $n$th infinitesimal
neighbourhood of $Z$ and denote $Z_{\mu, n} \subset V_\mu$
the $n$th infinitesimal neighbourhood of $Z_\mu$.
By the theorem on formal functions
(Cohomology of Spaces, Theorem
\ref{spaces-cohomology-theorem-formal-functions})
we have
$$
C^\wedge = \lim_n H^0(V' \times_V Z_n, \mathcal{O})
\quad\text{and}\quad
C_\mu^\wedge =
\lim_n H^0(V'_\mu \times_{V_\mu} Z_{\mu, n}, \mathcal{O})
$$
where $C^\wedge$ and $C_\mu^\wedge$ are the completion with
respect to $I$ and $I_\mu$.
Combining the completion of the morphism
$x'_\mu : V'_\mu \to X'$ with the morphism $g : X'_{/T'} \to W$ we obtain
$$
g \circ x'_{\mu, /Z_\mu} :
V'_{\mu, /Z_\mu} = \colim V_\mu' \times_{V_\mu} Z_{\mu, n}
\longrightarrow
W
$$
and hence by the description of the completion
$C_\mu^\wedge$ above we obtain a continuous ring homomorphism
$$
(g \circ x'_{\mu, /Z_\mu})^\sharp : B \longrightarrow C_\mu^\wedge
$$
The fact that $V' \to V$, $\hat x'$, $x'$ witnesses the compatibility
between $u'$ and $\hat x$ implies the
commutativity of the following diagram
$$
\xymatrix{
C_\mu^\wedge \ar[r] &
C^\wedge \\
B \ar[u]^{(g \circ x'_{\mu, /Z_\mu})^\sharp} \ar[r]^{(\hat x)^\sharp} &
A^\wedge \ar[u]
}
$$

\medskip\noindent
Proof of surjectivity: more algebra arguments. Recall that the finite
$A$-modules $\Ker(A \to C)$ and $\Coker(A \to C)$ are annihilated
by a power of $I$ and similarly the finite $A_\mu$-modules
$\Ker(A_\mu \to C_\mu)$ and $\Coker(A_\mu \to C_\mu)$ are annihilated
by a power of $I_\mu$. This implies that these modules are
equal to their completions. Since $I$-adic completion on the category of
finite $A$-modules is exact (see
Algebra, Section \ref{algebra-section-completion-noetherian})
it follows that we have
$$
\Coker(A^\wedge \to C^\wedge) = \Coker(A \to C)
$$
and similarly for kernels and for the maps $A_\mu \to C_\mu$.
Of course we also have
$$
\Ker(A \to C) = \colim \Ker(A_\mu \to C_\mu)
\quad\text{and}\quad
\Coker(A \to C) = \colim \Coker(A_\mu \to C_\mu)
$$
Recall that $S = \Spec(R)$ is affine. All of the ring maps
above are $R$-algebra homomorphisms as all of the morphisms
are morphisms over $S$. By
Algebraization of Formal Spaces, Lemma
\ref{restricted-lemma-Noetherian-finite-type-red}
we see that $B$ is topologically of finite type over $R$.
Say $B$ is topologically generated by $b_1, \ldots, b_n$.
Pick some $\mu$ (for example $\lambda$) and consider
the elements
$$
\text{images of }
(g \circ x'_{\mu, /Z_\mu})^\sharp(b_1)
, \ldots,
(g \circ x'_{\mu, /Z_\mu})^\sharp(b_n)
\text{ in }\Coker(A_\mu \to C_\mu)
$$
The image of these elements in $\Coker(\alpha)$ are zero
by the commutativity of the square above. Since
$\Coker(A \to C) = \colim \Coker(A_\mu \to C_\mu)$ and these
cokernels are equal to their completions
we see that after increasing $\mu$ we may assume these
images are all zero. This means that the continuous
homomorphism $(g \circ x'_{\mu, /Z_\mu})^\sharp$ has image contained
in $\Im(A_\mu \to C_\mu)$.
Choose elements $a_{\mu, j} \in (A_\mu)^\wedge$ mapping to
$(g \circ x'_{\mu, /Z_\mu})^\sharp(b_1)$ in $(C_\mu)^\wedge$.
Then $a_{\mu, j} \in A_\mu^\wedge$ and $(\hat x)^\sharp(b_j) \in A^\wedge$
map to the same element of $C^\wedge$ by the commutativity of the
square above. Since
$\Ker(A \to C) = \colim \Ker(A_\mu \to C_\mu)$ and these kernels
are equal to their completions, we may after increasing
$\mu$ adjust our choices of $a_{\mu, j}$ such that
the image of $a_{\mu, j}$ in $A^\wedge$ is equal to $(\hat x)^\sharp(b_j)$.

\medskip\noindent
Proof of surjectivity: final algebra arguments.
Let $\mathfrak b \subset B$ be the ideal of topologically nilpotent
elements. Let $J \subset R[x_1, \ldots, x_n]$ be the ideal
consisting of those $h(x_1, \ldots, x_n)$ such that
$h(b_1, \ldots, b_n) \in \mathfrak b$. Then we get a continuous
surjection of topological $R$-algebras
$$
\Phi : R[x_1, \ldots, x_n]^\wedge \longrightarrow B,\quad
x_j \longmapsto b_j
$$
where the completion on the left hand side is with respect to $J$.
Since $R[x_1, \ldots, x_n]$ is Noetherian we can choose
generators $h_1, \ldots, h_m$ for $J$. By the commutativity
of the square above we see that $h_j(a_{\mu, 1}, \ldots, a_{\mu, n})$ is
an element of $A_\mu^\wedge$ whose image in $A^\wedge$ is contained
in $IA^\wedge$. Namely, the ring map $(\hat x)^\sharp$ is continuous
and $IA^\wedge$ is the ideal of topological nilpotent elements
of $A^\wedge$ because $A^\wedge/IA^\wedge = A/I$ is reduced.
(See Algebra, Section \ref{algebra-section-completion-noetherian}
for results on completion in Noetherian rings.)
Since $A/I = \colim A_\mu/I_\mu$ we conclude that after increasing
$\mu$ we may assume $h_j(a_{\mu, 1}, \ldots, a_{\mu, n})$ is in
$I_\mu A_\mu^\wedge$. In particular the elements
$h_j(a_{\mu, 1}, \ldots, a_{\mu, n})$ of $A_\mu^\wedge$
are topologically nilpotent in $A_\mu^\wedge$.
Thus we obtain a continuous $R$-algebra homomorphism
$$
\Psi : R[x_1, \ldots, x_n]^\wedge \longrightarrow A_\mu^\wedge,\quad
x_j \longmapsto a_{\mu, j}
$$
In order to conclude what we want, we need to see if $\Ker(\Phi)$ is
annihilated by $\Psi$. This may not be true, but we can achieve
this after increasing $\mu$. Indeed, since
$R[x_1, \ldots, x_n]^\wedge$ is Noetherian,
we can choose generators $g_1, \ldots, g_l$ of the ideal
$\Ker(\Phi)$. Then we see that
$$
\Psi(g_1), \ldots, \Psi(g_l) \in
\Ker(A_\mu^\wedge \to C_\mu^\wedge) = \Ker(A_\mu \to C_\mu)
$$
map to zero in $\Ker(A \to C) = \colim \Ker(A_\mu \to C_\mu)$.
Hence increasing $\mu$ as before we get the desired result.

\medskip\noindent
Proof of surjectivity: mopping up. The continuous ring homomorphism
$B \to (A_\mu)^\wedge$ constructed above determines a morphism
$\hat x_\mu : V_{\mu, /Z_\mu} \to W$.
The compatibility of $\hat x_\mu$ and $u'_\mu$ follows
from the fact that the ring map $B \to (A_\mu)^\wedge$
is by construction compatible with the ring map $A_\mu \to C_\mu$.
In fact, the compatibility will be witnessed by the proper morphism
$V'_\mu \to V_\mu$ and the morphisms
$x'_\mu$ and $\hat x'_\mu = x'_{\mu, /Z_\mu}$ we used in
the construction. This finishes the proof.
\end{proof}
